\documentclass[10pt,a4paper]{amsart}
\usepackage{amsmath,amssymb,amsthm,graphicx}
% Retained minimal preamble for the Stacks Project chapter sites-cohomology.tex (Cohomology on Sites).
% Verbatim excerpt of the chapter preamble: Xy-pic setup (native preamble lines 14-18) and the native
% theorem declarations with the mathematical macros (native preamble lines 154-217). Copied mechanically;
% nothing added, removed or redefined. The wrapper class and the amsmath/amssymb/amsthm/graphicx base
% packages are supplied by the assembler, so the class/externaldocument part of the native preamble is omitted.
% For commutative diagrams we use Xy-pic
\usepackage[all]{xy}

% We use 2cell for 2-commutative diagrams.
\xyoption{2cell}
% Theorem environments.
%
\theoremstyle{plain}
\newtheorem{theorem}[subsection]{Theorem}
\newtheorem{proposition}[subsection]{Proposition}
\newtheorem{lemma}[subsection]{Lemma}

\theoremstyle{definition}
\newtheorem{definition}[subsection]{Definition}
\newtheorem{example}[subsection]{Example}
\newtheorem{exercise}[subsection]{Exercise}
\newtheorem{situation}[subsection]{Situation}

\theoremstyle{remark}
\newtheorem{remark}[subsection]{Remark}
\newtheorem{remarks}[subsection]{Remarks}

\numberwithin{equation}{subsection}

% Macros
%
\def\lim{\mathop{\mathrm{lim}}\nolimits}
\def\colim{\mathop{\mathrm{colim}}\nolimits}
\def\Spec{\mathop{\mathrm{Spec}}}
\def\Hom{\mathop{\mathrm{Hom}}\nolimits}
\def\Ext{\mathop{\mathrm{Ext}}\nolimits}
\def\SheafHom{\mathop{\mathcal{H}\!\mathit{om}}\nolimits}
\def\SheafExt{\mathop{\mathcal{E}\!\mathit{xt}}\nolimits}
\def\Sch{\mathit{Sch}}
\def\Mor{\mathop{\mathrm{Mor}}\nolimits}
\def\Ob{\mathop{\mathrm{Ob}}\nolimits}
\def\Sh{\mathop{\mathit{Sh}}\nolimits}
\def\NL{\mathop{N\!L}\nolimits}
\def\CH{\mathop{\mathrm{CH}}\nolimits}
\def\proetale{{pro\text{-}\acute{e}tale}}
\def\etale{{\acute{e}tale}}
\def\QCoh{\mathit{QCoh}}
\def\Ker{\mathop{\mathrm{Ker}}}
\def\Im{\mathop{\mathrm{Im}}}
\def\Coker{\mathop{\mathrm{Coker}}}
\def\Coim{\mathop{\mathrm{Coim}}}

% Boxtimes
%
\DeclareMathSymbol{\boxtimes}{\mathbin}{AMSa}{"02}

%
% Macros for moduli stacks/spaces
%
\def\QCohstack{\mathcal{QC}\!\mathit{oh}}
\def\Cohstack{\mathcal{C}\!\mathit{oh}}
\def\Spacesstack{\mathcal{S}\!\mathit{paces}}
\def\Quotfunctor{\mathrm{Quot}}
\def\Hilbfunctor{\mathrm{Hilb}}
\def\Curvesstack{\mathcal{C}\!\mathit{urves}}
\def\Polarizedstack{\mathcal{P}\!\mathit{olarized}}
\def\Complexesstack{\mathcal{C}\!\mathit{omplexes}}
% \Pic is the operator that assigns to X its picard group, usage \Pic(X)
% \Picardstack_{X/B} denotes the Picard stack of X over B
% \Picardfunctor_{X/B} denotes the Picard functor of X over B
\def\Pic{\mathop{\mathrm{Pic}}\nolimits}
\def\Picardstack{\mathcal{P}\!\mathit{ic}}
\def\Picardfunctor{\mathrm{Pic}}
\def\Deformationcategory{\mathcal{D}\!\mathit{ef}}
\usepackage{cleveref}
\crefname{originallemma}{original lemma}{original lemmas}
\Crefname{originallemma}{Original lemma}{Original lemmas}
\makeatletter
\newlabel{homology-lemma-characterize-weak-serre-subcategory}{{0754}{source}{}{sourcepointer-1}{}}
\newlabel{homology-lemma-characterize-weak-serre-subcategory@cref}{{[originallemma][0][]0754}{[1][0][]source}{}{}{}}
\makeatother
\crefname{originallemma}{original lemma}{original lemmas}
\Crefname{originallemma}{Original lemma}{Original lemmas}
\makeatletter
\newlabel{lemma-V-C-all-n-general}{{0EZF}{source}{}{sourcepointer-2}{}}
\newlabel{lemma-V-C-all-n-general@cref}{{[originallemma][0][]0EZF}{[1][0][]source}{}{}{}}
\makeatother
\crefname{originallemma}{original lemma}{original lemmas}
\Crefname{originallemma}{Original lemma}{Original lemmas}
\makeatletter
\newlabel{lemma-equivalence-bounded}{{0D7U}{source}{}{sourcepointer-3}{}}
\newlabel{lemma-equivalence-bounded@cref}{{[originallemma][0][]0D7U}{[1][0][]source}{}{}{}}
\makeatother
\crefname{originalreference}{original reference}{original references}
\Crefname{originalreference}{Original reference}{Original references}
\makeatletter
\newlabel{situation-compare}{{0EZ3}{source}{}{sourcepointer-4}{}}
\newlabel{situation-compare@cref}{{[originalreference][0][]0EZ3}{[1][0][]source}{}{}{}}
\makeatother
\title{The category of sheaves of the coarser topology is a weak Serre subcategory and R epsilon\_* is an equivalence on bounded below complexes}
\author{The Stacks Project Authors}
\date{stacks-project:sites-cohomology.tex:local-snapshot-20261011; GNU FDL 1.2 or later}
\begin{document}
\maketitle
\noindent\textit{Editorial note.} All delivered mathematical text is the author's own, copied from the pinned Stacks Project chapter file sites-cohomology.tex; no statement or proof step was rewritten, reordered or removed.\par\medskip
\noindent\textit{Editorial note.} The counted proof is the author's complete argument: the proof is written out in full: the weak Serre axioms for A\_X are verified by transporting kernels, cokernels and extensions along epsilon\_\{X,*\} and epsilon\_X\textasciicircum{}\{-1\}, and the equivalence is obtained from the equivalence criterion for bounded below derived categories once the conditions (V\_n) for all n are read off from the V-C lemma.\par\medskip
\noindent\textit{Editorial note.} The label lemma-cohomological-descent-general names the counted statement, and the proof environment immediately following it in the source is delivered verbatim as its proof.\par\medskip
\noindent\textit{Editorial note.} The statement refers to the standing situation of the comparison of two topologies, which is named in the statement and kept as an external pointer with its canonical Stacks tag.\par\medskip
\noindent\textit{Editorial note.} The preamble is the chapter's own retained minimal preamble (sources/stacks-sites-cohomology/retained-preamble.tex), a verbatim excerpt of the native preamble: the Xy-pic setup and the native theorem declarations with the mathematical macros.\par\medskip
\noindent\textit{Allowed standard prerequisites (editorial).}\begin{itemize}
\item Situation of two topologies on one category and the morphism epsilon. Same chapter, the situation comparing a topology tau' with a coarser topology tau on C and the morphism of sites epsilon : C\_tau -> C\_tau'. Supplies the categories of sheaves, the functors epsilon\_*, epsilon\textasciicircum{}\{-1\} and the subcategories A\_X, A'\_X appearing in the statement.
\item Characterisation of weak Serre subcategories. Homology, the lemma listing the exactness conditions characterising a weak Serre subcategory. Supplies the four conditions that are verified for A\_X at the start of the proof.
\item Equivalence of bounded below derived categories from an equivalence on the hearts. Same chapter, the lemma deducing an equivalence D\textasciicircum{}+\_\{A'\} -> D\textasciicircum{}+\_A from an equivalence of A' with A and the vanishing of the derived functors on A'. Supplies the final step of the proof, after the conditions (V\_n) for all n have been obtained.
\end{itemize}\medskip
\ifdefined\hypertarget\hypertarget{sourcepointer-1}{}\fi
\noindent\textit{Source reference.} Original lemma 0754: sources/stacks-chow/tags.txt, source lines 9118--9118. Prerequisite pointer invoked at proof line(s) 7377; it names a standard fact of a sibling chapter that the counted argument uses, and every step of the counted proof built on it is retained verbatim in the delivered body.\par\medskip
\ifdefined\hypertarget\hypertarget{sourcepointer-2}{}\fi
\noindent\textit{Source reference.} Original lemma 0EZF: sources/stacks-sites-cohomology/pinned-original-source.tex, source lines 7292--7292. Prerequisite pointer invoked at proof line(s) 7406; it names a standard fact of the same chapter that the counted argument uses, and every step of the counted proof built on it is retained verbatim in the delivered body.\par\medskip
\ifdefined\hypertarget\hypertarget{sourcepointer-3}{}\fi
\noindent\textit{Source reference.} Original lemma 0D7U: sources/stacks-sites-cohomology/pinned-original-source.tex, source lines 6391--6391. Prerequisite pointer invoked at proof line(s) 7421; it names a standard fact of the same chapter that the counted argument uses, and every step of the counted proof built on it is retained verbatim in the delivered body.\par\medskip
\ifdefined\hypertarget\hypertarget{sourcepointer-4}{}\fi
\noindent\textit{Source reference.} Original reference 0EZ3: sources/stacks-sites-cohomology/pinned-original-source.tex, source lines 6841--6841. Prerequisite pointer invoked at proof line(s) 7362; it names a standard fact of the same chapter that the counted argument uses, and every step of the counted proof built on it is retained verbatim in the delivered body.\par\medskip
\begin{lemma}
\label{lemma-cohomological-descent-general}
In Situation \ref{situation-compare}. For any $X$ in
$\mathcal{C}$ the category
$\mathcal{A}_X \subset \textit{Ab}(\mathcal{C}_\tau/X)$
is a weak Serre subcategory and the functor
$$
R\epsilon_{X, *} :
D^+_{\mathcal{A}_X}(\mathcal{C}_\tau/X)
\longrightarrow
D^+_{\mathcal{A}'_X}(\mathcal{C}_{\tau'}/X)
$$
is an equivalence with quasi-inverse given by $\epsilon_X^{-1}$.
\end{lemma}

\begin{proof}
We need to check the conditions listed in
Homology, Lemma \ref{homology-lemma-characterize-weak-serre-subcategory}
for $\mathcal{A}_X$. If $\varphi : \mathcal{F} \to \mathcal{G}$ is a map
in $\mathcal{A}_X$, then $\epsilon_{X, *}\varphi :
\epsilon_{X, *}\mathcal{F} \to \epsilon_{X, *}\mathcal{G}$
is a map in $\mathcal{A}'_X$. Hence $\Ker(\epsilon_{X, *}\varphi)$ and
$\Coker(\epsilon_{X, *}\varphi)$ are objects of $\mathcal{A}'_X$
as this is a weak Serre subcategory of $\textit{Ab}(\mathcal{C}_{\tau'}/X)$.
Applying $\epsilon_X^{-1}$ we obtain an exact sequence
$$
0 \to
\epsilon_X^{-1}\Ker(\epsilon_{X, *}\varphi) \to
\mathcal{F} \to \mathcal{G} \to
\epsilon_X^{-1}\Coker(\epsilon_{X, *}\varphi) \to 0
$$
and we see that $\Ker(\varphi)$ and $\Coker(\varphi)$ are in
$\mathcal{A}_X$. Finally, suppose that
$$
0 \to \mathcal{F}_1 \to \mathcal{F}_2 \to \mathcal{F}_3 \to 0
$$
is a short exact sequence in $\textit{Ab}(\mathcal{C}_\tau/X)$
with $\mathcal{F}_1$ and $\mathcal{F}_3$ in $\mathcal{A}_X$.
Then applying $\epsilon_{X, *}$ we obtain an exact sequence
$$
0 \to 
\epsilon_{X, *}\mathcal{F}_1 \to
\epsilon_{X, *}\mathcal{F}_2 \to
\epsilon_{X, *}\mathcal{F}_3 \to
R^1\epsilon_{X, *}\mathcal{F}_1 = 0
$$
Vanishing by Lemma \ref{lemma-V-C-all-n-general}.
Hence $\epsilon_{X, *}\mathcal{F}_2$ is in $\mathcal{A}'_X$
as this is a weak Serre subcategory of $\textit{Ab}(\mathcal{C}_{\tau'}/X)$.
Pulling back by $\epsilon_X$ we conclude that
$\mathcal{F}_2$ is in $\mathcal{A}_X$.

\medskip\noindent
Thus $\mathcal{A}_X$ is a weak Serre subcategory of
$\textit{Ab}(\mathcal{C}_\tau/X)$ and it makes sense
to consider the category $D^+_{\mathcal{A}_X}(\mathcal{C}_\tau/X)$.
Observe that $\epsilon_X^{-1} : \mathcal{A}'_X \to \mathcal{A}_X$
is an equivalence and that
$\mathcal{F}' \to R\epsilon_{X, *}\epsilon_X^{-1}\mathcal{F}'$
is an isomorphism for $\mathcal{F}'$ in $\mathcal{A}'_X$ since we
have $(V_n)$ for all $n$ by Lemma \ref{lemma-V-C-all-n-general}.
Thus we conclude by Lemma \ref{lemma-equivalence-bounded}.
\end{proof}

\end{document}
